\chapter{Терминологический реестр}
\label{app:terminology}

\section{Назначение реестра}

Глоссарий репозитория является нормативным источником терминов для
документации, протоколов, отчётов и текста диссертации. Поэтому терминология
работы не должна существовать в двух несвязанных вариантах: отдельно в
\code{docs/glossary.md} и отдельно в основном тексте. Настоящее приложение
включает полный русский реестр всех стабильных идентификаторов \code{TERM-*},
которые существуют в нормативном глоссарии на момент сборки диссертации.

Ключевые понятия, непосредственно необходимые для постановки исследовательских
вопросов и интерпретации результатов, вводятся в основных главах. В частности,
глава~\ref{ch:related-work} раскрывает PC-TREF, PC-CATM и QWake-PC/QWake-FP до
описания экспериментов. Приложение выполняет другую функцию: оно гарантирует,
что ни один нормативный термин не остаётся только внутренним понятием
репозитория и что читатель может установить его каноническое значение и границу
употребления внутри самой диссертации.

Наличие термина в этом реестре не означает эмпирического подтверждения
соответствующего механизма. Разделы «Перспективные расширения теории и
управления» и «Зарезервированные обобщения вычислительного действия» содержат
понятия, смысл которых зафиксирован для однозначности, но которые находятся за
границей обязательных эмпирических утверждений текущей работы.

\section{Полнота и источник}

Реестр ниже формируется автоматически из \code{docs/glossary.md} во время
\code{make thesis}. Сборка проверяет совпадение упорядоченного множества
идентификаторов с \code{docs/glossary_EN.md}, уникальность идентификаторов,
наличие русского определения и правила употребления для каждого термина и
полноту включения в диссертацию. На текущем состоянии репозитория ожидается
117 из 117 терминов.

\input{generated/terminology_registry}
